% part: intuitionistic-logic
% chapter: introduction
% section: axiomatic-derivations

\documentclass[../../../include/open-logic-section]{subfiles}

\begin{document}

\olfileid[pt]{int}{int}{axd}

\olsection{\usetoken{P}{derivation} axiomáticas}

As !!{derivation}s axiomáticas da lógica proposicional intuicionista
são os sistemas de !!{derivation} conceitualmente mais simples e,
historicamente, os primeiros. Elas funcionam exatamente como na lógica
proposicional clássica.

\begin{defn}[!!^{derivability}]
Se $\Gamma$ é um conjunto de !!{formula}s de $\Lang L$, então uma
\emph{!!{derivation}} a partir de $\Gamma$ é uma sequência finita
$!A_1$, \dots,~$!A_n$ de !!{formula}s tal que, para cada $i \le n$,
vale uma das seguintes condições:
\begin{enumerate}
\item $!A_i \in \Gamma$; ou
\item $!A_i$ é um axioma; ou
\item $!A_i$ decorre de algum $!A_j$ e de algum $!A_k$, com $j < i$ e
  $k < i$, por modus ponens, isto é, $!A_k \ident !A_j \lif !A_i$.
\end{enumerate}
\end{defn}

\begin{defn}[Axiomas]
O conjunto $\PAx$ de \emph{axiomas} da lógica proposicional intuicionista
consiste em todas as !!{formula}s das seguintes formas:
\begin{align}
  & (!A \land !B) \lif !A \ollabel{ax:land1}\\
  & (!A \land !B) \lif !B \ollabel{ax:land2}\\
  & !A \lif (!B \lif (!A \land !B)) \ollabel{ax:land3}\\
  & !A \lif (!A \lor !B) \ollabel{ax:lor1}\\
  & !A \lif (!B \lor !A) \ollabel{ax:lor2}\\
  & (!A \lif !C) \lif ((!B \lif !C) \lif ((!A \lor !B) \lif !C)) \ollabel{ax:lor3}\\
  & !A \lif (!B \lif !A) \ollabel{ax:lif1}\\
  & (!A \lif (!B \lif !C)) \lif ((!A \lif !B) \lif (!A \lif !C)) \ollabel{ax:lif2}\\
  & \lfalse \lif !A \ollabel{ax:lfalse1}
\end{align}
\end{defn}

\begin{defn}[!!^{derivability}]
!!^a{formula}~$!A$ é \emph{!!{derivable}} a partir de $\Gamma$, o que
se escreve $\Gamma \Proves !A$, se existe !!a{derivation} a partir de
$\Gamma$ que termina em $!A$.
\end{defn}

\begin{defn}[Teoremas]
!!^a{formula}~$!A$ é um \emph{teorema} se existe !!a{derivation}
de~$!A$ a partir do conjunto vazio. Escrevemos $\Proves !A$ se $!A$ é
um teorema e $\Proves/ !A$ se não é.
\end{defn}

\begin{prop}
  Se $\Gamma \Proves !A$ no sistema axiomático de !!{derivation} da
  lógica intuicionista, então $\Gamma \Proves !A$ também na lógica
  clássica. Em particular, se $!A$ é um teorema intuicionista, também é
  um teorema clássico.
\end{prop}

\begin{proof}
  Todo axioma intuicionista também é um axioma clássico; portanto, toda
  !!{derivation} na lógica intuicionista também é !!a{derivation} na
  lógica clássica.
\end{proof}

\end{document}
